% Part: counterfactuals
% Chapter: minimal-change-semantics
% Section: transitivity

\documentclass[../../../include/open-logic-section]{subfiles}

\begin{document}

\olfileid{cnt}{min}{tra}

\olsection{సంక్రామకత్వం}

భౌతిక సోపాధికానికి గొలుసు నియమం
వర్తిస్తుంది: $!A \lif !B, !B
\lif !C \Entails !A \lif !C$.
మరో మాటలో, భౌతిక సోపాధికం
సంక్రామకం. ప్రతివాస్తవాలకు కూడా
ఇది వర్తిస్తుందా? స్టాల్నేకర్
ఇచ్చిన కింది ఉదాహరణను పరిశీలించండి.
\begin{quote}
  జే.~ఎడ్గర్ హూవర్ రష్యన్‌గా పుట్టి
  ఉంటే, కమ్యూనిస్టు అయ్యేవారు.

  జే.~ఎడ్గర్ హూవర్ కమ్యూనిస్టు అయి
  ఉంటే, దేశద్రోహి అయ్యేవారు.

  కాబట్టి జే.~ఎడ్గర్ హూవర్ రష్యన్‌గా
  పుట్టి ఉంటే, దేశద్రోహి అయ్యేవారు.
\end{quote}
హూవర్ వాస్తవంలో పుట్టిన అదే
సమయంలో అమెరికాలో కాకుండా
రష్యాలో పుట్టి ఉంటే, సోవియట్
యూనియన్‌లో పెరిగి కమ్యూనిస్టు
అయ్యేవారని (అనుకుందాం). అందువల్ల
మొదటి పూర్వాపేక్ష సత్యం.
అలాగే రెండో పూర్వాపేక్షను
విడిగా పరిశీలిస్తే అదీ సత్యమే.
కానీ ముగింపు అసత్యం: హూవర్
సోవియట్ యూనియన్‌లో పుట్టి
ఉంటే బహుశా ఉత్సాహపూరిత
కమ్యూనిస్టుగా ఉండేవారు,
తన దేశానికి దేశద్రోహి అయ్యేవారు కాదు.
ఈ సహజ సత్యమూల్య నియామకానికి
స్టాల్నేకర్--లూయిస్ వివరణ
మద్దతిస్తుంది. హూవర్ జన్మస్థలం
మాత్రమే మారిన మన లోకానికి
అత్యంత సమీప సాధ్యలోకంలో,
హూవర్ సోవియట్ యూనియన్‌కు
మంచి పౌరుడిగా పెరుగుతారు.
మొదటి పూర్వాపేక్ష పూర్వపక్షం,
ముగింపు పూర్వపక్షం రెండూ
సత్యమయ్యే అత్యంత సమీప లోకం
ఇదే; ఆ లోకంలో హూవర్
కమ్యూనిస్టు పార్టీకి విధేయుడు,
అందువల్ల దేశద్రోహి కాదు.
అయితే రెండో పూర్వాపేక్షను
మూల్యాంకనం చేయడానికి వేరే
లోకాన్ని చూడాలి: హూవర్
కమ్యూనిస్టుగా ఉండే అత్యంత
సమీప లోకం. అందులో ఆయన
అమెరికాలో పుట్టి,
కమ్యూనిస్టుగా మారి,
తద్వారా దేశద్రోహి అయ్యారు.\footnote{ఈ
  ఉదాహరణ బలాన్ని అర్థం చేసుకోవాలంటే
  కొన్ని అధిభౌతిక, రాజకీయ
  పరికల్పనలు స్వీకరించాలి: ఉదాహరణకు,
  హూవర్ రష్యన్ తల్లిదండ్రులకు
  పుట్టి ఉండడం సాధ్యమని లేదా
  1950ల అమెరికాలోని కమ్యూనిస్టులు
  తమ దేశానికి దేశద్రోహులని.}

\begin{prob}
  ప్రతివాస్తవాల సంక్రామకత్వం
  విఫలమయ్యే నమ్మదగిన సహజ
  ఉదాహరణను కనుగొనండి.
\end{prob}

\begin{ex}\ollabel{ex:trans-counterex}
  గోళ అర్థవిచారంలో ఈ నిగమనం
  చెల్లదు; అంటే $p
  \cif q, q \cif r \Entails/ p \cif r$.
  $W = \{w, w_1, w_2\}$,
  $O_w = \{\{w\}, \{w,
  w_1\}, \{w, w_1, w_2\}\}$,
  $V(p) = \{w_2\}$,
  $V(q) = \{w_1, w_2\}$,
  $V(r) = \{w_1\}$
  గల నమూనా $\mModel{M}
  = \tuple{W, O, V}$ను
  పరిశీలించండి. $p$ సత్యమయ్యే
  లోకాన్ని కలిగిన గోళం
  $S = \{w, w_1,
  w_2\}$ ఉంది; దాని
  అన్ని లోకాల వద్ద
  $p \lif q$ సత్యం.
  కాబట్టి $\mSat{M}{p
    \cif q}[w]$. అలాగే
  $q$ సత్యమయ్యే లోకాన్ని
  కలిగిన గోళం $S' = \{w,
  w_1\}$ ఉంది; దాని
  అన్ని లోకాల వద్ద
  $\mSat{M}{q \lif r}$,
  కాబట్టి $\mSat{M}{q \cif r}[w]$.
  \sourcecorrection{OLTECNTTRA-001}{మూల నమూనా వివరణలో $q \lif r$ అన్ని లోకాల వద్ద సత్యమని గద్యంగా చెప్పి, సంకేతంలో మాత్రం అసత్యమని వ్రాసింది. ఇచ్చిన $V(q)$, $V(r)$ విలువలకు అనుగుణంగా సత్య సంకేతానికి సరిచేశాం.}
  అయితే $p$ సత్యమయ్యే
  గోళం $\{w,
  w_1, w_2\}$లో
  ఒక లోకం, అంటే~$w_2$,
  వద్ద $\mSat/{M}{p \lif
    r}[w_2]$.
\end{ex}

\begin{prob}
\olref[cnt][min][tra]{ex:trans-counterex}లోని
ప్రతిదృష్టాంతానికి సరిపోయే
గోళ చిత్రాన్ని గీయండి.
\end{prob}

\begin{prob}
  \olref[cnt][min][tra]{ex:trans-counterex}లో
  లోకం $w_2$ వద్ద హూవర్
  రష్యాలో పుట్టి, కమ్యూనిస్టుగా
  ఉన్నా దేశద్రోహి కాదు;
  $w_1$ వద్ద ఆయన
  అమెరికాలో పుట్టి,
  కమ్యూనిస్టుగా ఉండి,
  దేశద్రోహి. ఈ నమూనాలో
  $w_1$,~$w$కు
  $w_2$కన్నా దగ్గర.
  ఇది తప్పనిసరా?
  హూవర్ కమ్యూనిస్టుగా
  ఉండే సాధ్యతకన్నా
  రష్యాలో పుట్టే సాధ్యత
  మరింత దూరమని అనుకోని
  ప్రతిదృష్టాంతం ఇవ్వగలరా?
\end{prob}

\end{document}
